\section{Design}
\label{sec:design}

\subsection{Architecture Overview}
\label{sec:design:arch}

\sys is a slab-based memory allocator with integrated compression
support.  It interposes on \texttt{malloc}, \texttt{free},
\texttt{realloc}, and related functions via the alloc8 interposition
framework, which uses \texttt{DYLD\_INSERT\_LIBRARIES} on macOS and
\texttt{LD\_PRELOAD} on Linux.  Applications link against the system
allocator normally; \sys replaces it at load time.

Figure~\ref{fig:arch} shows the high-level architecture.  The
allocator manages heap pages through a virtual memory region (up to
16\,GiB).  Each page transitions through a state machine:

\begin{enumerate}
\item \textbf{EMPTY}: Page has never been allocated or has been
  returned to the OS.
\item \textbf{ACTIVE}: Page contains live objects with full
  read/write access.
\item \textbf{ACTIVE\_MONITORING}: Page has been marked read-only
  (\texttt{mprotect} with \texttt{PROT\_READ}) to detect future writes.
\item \textbf{COMPRESSING}: Compressor is actively compressing
  this page.
\item \textbf{COMPRESSED}: Page data is stored in compressed form;
  the physical page has been decommitted.
\end{enumerate}

\begin{figure}[t]
\centering
\small
\begin{verbatim}
 EMPTY --> ACTIVE --> ACTIVE_MONITORING
             ^              |
             |              v
             +-------- COMPRESSING
             |              |
             |              v
             +-------- COMPRESSED
\end{verbatim}
\caption{Page state machine.  Transitions from COMPRESSED or
  ACTIVE\_MONITORING back to ACTIVE are triggered by the signal
  handler on page fault.  The COMPRESSING $\to$ ACTIVE transition
  occurs when a fault preempts an in-progress compression.}
\label{fig:arch}
\end{figure}

A background compressor thread wakes every second
(\texttt{kCompressIntervalMs} = 1000\,ms) and executes three phases:
(1)~update per-page cold counts based on access bits,
(2)~compress pages whose cold count exceeds a threshold, and
(3)~set up monitoring (\texttt{PROT\_READ}) on remaining active pages.

When the application accesses a compressed or monitored page, a
SIGSEGV or SIGBUS signal fires.  The fault handler decompresses the
page (or restores write permission for monitored pages) and returns,
allowing the faulting instruction to retry transparently.

\subsection{Call-Site Arena Segregation}
\label{sec:design:arenas}

\sys maintains four arenas (\texttt{kNumArenas} = 4, configurable as a
power of two).  Each arena has independent slab free lists for all 36
size classes, organized as a flat array
\texttt{slabs\_[kNumArenas $\times$ kNumClasses]}.

On each \texttt{malloc} call, \sys selects an arena by hashing the
return address:

\begin{lstlisting}
uintptr_t h =
    __builtin_return_address(0) ^
    __builtin_return_address(1);
h ^= h >> 16;
arena = h & (kNumArenas - 1);
\end{lstlisting}

XORing two stack depths ensures different arena routing even when
allocations pass through a common wrapper function.  The hash maps
each unique (caller, grandcaller) pair to a deterministic arena.
Allocations from the same call site consistently land in the same
arena, producing pages with structurally homogeneous content.

On \texttt{free}, the thread cache drains objects back to their
originating arena by looking up each pointer's span in the page map
and reading the span's \texttt{arena\_id} field.  This prevents arena
mixing during deallocation.

\subsection{Adaptive Multi-Algorithm Compression}
\label{sec:design:adaptive}

\sys uses a two-tier compression strategy based on how long a page has
been cold:

\begin{itemize}
\item \textbf{Cold} ($\geq$ \texttt{kColdTicks} = 2 ticks, i.e.,
  2 seconds): Compress with LZ4 at acceleration level~1.  LZ4
  compresses at $\sim$1\,GB/s and decompresses at $\sim$4--6\,GB/s,
  making it suitable for pages that may be accessed again soon.

\item \textbf{Very cold} ($\geq$ \texttt{kVeryColdTicks} = 60 ticks,
  i.e., 1 minute): Recompress with zstd at level~9.  zstd achieves
  2--15 percentage points better compression than LZ4 on structured
  data (Table~\ref{tab:algo_compare}), at the cost of 10$\times$
  slower compression.  Decompression speed is comparable across
  levels (1.2--1.8\,GB/s).
\end{itemize}

\paragraph{Per-size-class skip statistics.}
Not all size classes compress well.  \sys maintains a 64-entry sliding
window of compression ratios for each size class.  The window records
the compression ratio (0--255, mapping to 0--100\%) of each attempted
compression.  If the average ratio across at least 32 samples falls
below 5\% space savings, \sys skips LZ4 compression attempts for that
size class.  The sliding window allows recovery: if the data pattern
changes and new pages compress well, the old poor-ratio samples age
out.

\subsection{Zero-on-Free}
\label{sec:design:zero}

When an object is freed, \sys replaces its contents with zeros.  This
optimization targets partially-filled slab pages, where freed slots
would otherwise contain stale data from their previous occupant.
Zeroed regions form long runs that LZ4 and zstd compress extremely
efficiently.

\sys uses a two-tier zeroing strategy to minimize the performance
impact on \texttt{free}:

\begin{itemize}
\item \textbf{Eager zeroing} ($\leq$128\,B): Small objects are zeroed
  with \texttt{memset} directly in \texttt{free()}.  The 128-byte
  threshold (\texttt{kZeroOnFreeMaxSize}) was chosen because objects
  this small are typically already in the L1 cache at free time, so
  zeroing adds negligible latency ($\sim$2\,ns on our test hardware).

\item \textbf{Deferred zeroing} ($>$128\,B): The compressor thread
  zeros freed slots in its scratch buffer after copying the page data
  but before compressing.  Zeroing uses non-temporal stores
  (\texttt{\_\_builtin\_nontemporal\_store}, compiled to ARM64
  \texttt{stnp} instructions) to avoid polluting the data cache with
  cold page contents.
\end{itemize}

The deferred strategy is critical for large objects.  Unconditional
zeroing of a 16\,KiB object in \texttt{free()} would add
$\sim$1.5\,$\mu$s of latency, a 70$\times$ increase over the baseline
21\,ns.  The threshold-based approach keeps \texttt{free()} latency
constant at $\sim$23\,ns across all size classes.

\subsection{Parallel Compressor}
\label{sec:design:parallel}

The compressor uses a coordinator-worker model with
\texttt{kCompressorWorkers} = 2 parallel workers (configurable).  The
coordinator thread wakes every second, partitions the committed page
range across workers, and dispatches each of the three phases.

\paragraph{Chunk bitmap.}  Rather than scanning all committed pages
linearly, the compressor maintains a bitmap with one bit per page in
64-page chunks.  A set bit indicates the page is non-EMPTY.  The
bitmap is rebuilt at each tick by scanning page states.  Iteration
uses \texttt{\_\_builtin\_ctzll} (count trailing zeros) to jump
directly to set bits, skipping large EMPTY regions in
$O(\text{active pages})$ rather than $O(\text{committed pages})$.

\paragraph{Per-worker state.}  Each worker has its own LZ4 state
buffer, zstd compression context (\texttt{ZSTD\_CCtx}), page-sized
scratch buffer, and compression output buffer.  Workers share only the
dictionary lookup table (read-only during compression) and the
sharded compressed data store.

\paragraph{Sharded CompressStore.}  Compressed page data is stored in
a pool with \texttt{kCompressStoreShards} = 8 independent shards,
each with its own spinlock and free lists.  Pages are assigned to
shards by \texttt{page\_idx \% 8}.  With 2 workers and 8 shards,
lock contention is near zero.

\subsection{Fault Handler}
\label{sec:design:fault}

\sys installs a signal handler for both SIGSEGV and SIGBUS.  The
handler receives the faulting address, looks up the page index in the
VM region, acquires the per-page lock, and takes action based on page
state:

\begin{itemize}
\item \textbf{COMPRESSED}: Recommit the physical page, decompress the
  stored data, release the compressed blob, transition to ACTIVE.
\item \textbf{COMPRESSING}: The compressor is mid-compression.
  Recommit the page (the compressor has a copy in its scratch buffer)
  and transition to ACTIVE.  The compressor detects the state change
  and abandons the compression.
\item \textbf{ACTIVE\_MONITORING}: A write to a read-only monitored
  page.  Restore \texttt{PROT\_READ|PROT\_WRITE}, mark the page as
  accessed, transition to ACTIVE.
\item \textbf{ACTIVE}: A race condition where Phase~3's batched
  \texttt{mprotect} (which sets \texttt{PROT\_READ} on a range of
  pages) overwrites a per-page \texttt{PROT\_RW} restoration from a
  concurrent fault handler.  Restore write permission.
\end{itemize}

\paragraph{Per-slot decompression contexts.}
The signal handler must not call \texttt{malloc}.  All decompression
state is pre-allocated from a bootstrap allocator.  \sys maintains 32
fault slots, each with a pre-allocated page buffer and an independent
\texttt{ZSTD\_DCtx}.  A faulting thread atomically acquires a slot via
compare-and-swap, decompresses into the slot's buffer, copies the
result to the faulting page, and releases the slot.  This design
supports up to 32 concurrent page faults without data races between
decompression contexts.

\paragraph{Prefetch.}
After decompressing a faulted page, the handler speculatively
decompresses up to \texttt{kPrefetchWindow} = 2 pages in each
direction, provided they belong to the same span and are in the
COMPRESSED state.  Prefetch uses non-blocking \texttt{tryLock} to
avoid deadlock with concurrent fault handlers.  Prefetch reduces
cold-access latency in sequential access patterns; memcached's cold
key re-access throughput improves by 54\% with prefetch enabled.

\paragraph{Bootstrap allocator.}
All \sys internal metadata (page state arrays, compressed data store,
decompression contexts, worker buffers) is allocated from a bump
allocator (\texttt{BootstrapAlloc}) backed by \texttt{mmap}.  The
bootstrap allocator never calls \texttt{malloc} and never frees memory.
This design is critical for the signal handler path, which must be
async-signal-safe.

\subsection{Syscall Compatibility}
\label{sec:design:syscall}

Kernel system calls access userspace buffers directly through
\texttt{copyin}/\texttt{copyout}.  If \sys has marked a buffer page
as \texttt{PROT\_READ} (monitoring) or \texttt{PROT\_NONE}
(compressed), the kernel returns EFAULT rather than triggering a
signal.  This affects any syscall that reads from or writes to
userspace memory: \texttt{read}, \texttt{write}, \texttt{recv},
\texttt{send}, \texttt{kevent}, and others.

\sys interposes on 20 system calls and C library I/O functions.
Before each call, the wrapper walks the buffer's pages, touching any
page managed by \sys to trigger decompression via the fault handler,
then pins the pages to prevent the compressor from re-protecting them
during the syscall.  After the syscall returns, pages are unpinned.

Page pins are implemented as per-page atomic counters
(\texttt{g\_page\_pins}).  The compressor skips pages with a nonzero
pin count in both Phase~2 (compression) and Phase~3 (monitoring).

\paragraph{Buffered I/O.}  On macOS, DYLD interposition only
intercepts cross-dylib calls.  When the application calls
\texttt{fread}, the C library internally calls \texttt{read} within
the same dylib (libSystem); \sys's \texttt{read} interposition does
not see this internal call.  \sys addresses this by also interposing
\texttt{fread}, \texttt{fgets}, \texttt{fgetc}, \texttt{fwrite}, and
\texttt{fflush}, warming both the user buffer and the \texttt{FILE}
struct's internal buffer (\texttt{stream->\_bf.\_base}).
Additionally, the compressor permanently pins the internal buffers of
\texttt{stdin}, \texttt{stdout}, and \texttt{stderr} on its first
tick, covering the intra-libSystem blind spot for standard streams.
